% Part: second-order-logic
% Chapter: metatheory
% Section: undecidability-and-axiomatizability

\documentclass[../../../include/open-logic-section]{subfiles}

\begin{document}

\olfileid{sol}{met}{nax}

\olsection{Logika Orde Kapindho Ora Bisa Diaksiomatisasi}


\begin{thm}
\ollabel{thm:sol-undecidable}
Logika orde kapindho ora bisa diputusake.
\end{thm}

\begin{proof}
!!{sentence} orde kapisan valid ing logika orde kapisan yen lan
mung yen valid ing logika orde kapindho, dene logika orde
kapisan ora bisa diputusake.
\end{proof}

\begin{thm}
\ollabel{cor:sol-not-axiomatizable}
Ora ana sistem !!{derivation} kang sahih lan jangkep kanggo
logika orde kapindho.
\end{thm}

\begin{proof}
Upamana $!A$ iku !!a{sentence} ing basa aritmetika.
$\Sat{N}{!A}$ yen lan mung yen $\Th{PA^2} \Entails !A$.
Upamana $!P$ iku konjungsi sangang aksioma saka~$\Th{PA^2}$.
$\Th{PA^2} \Entails !A$ yen lan mung yen
$\Entails !P \lif !A$, yaiku, kanggo saben struktur,
$\Sat{M}{!P \lif !A}$. Cathetan koreksi sumber OLPL-808: kurung kurawal kanggo pratelan kepuasan ing naskah sumber salah papan; wujud ing kene njaga maksud validitas implikasi. Saiki gatekna !!{sentence}
$\lforall[z][\lforall[u][\lforall[u'][\lforall[u''][\lforall[L][(!P' \lif !A')]]]]]$
kang diasilake kanthi ngganti $\Obj{0}$ dadi $z$,
$\prime$ dadi variabel fungsi siji-papan $u$,
$+$ lan $\times$ dadi variabel fungsi rong-papan
$u'$ lan $u''$, kanthi urut, sarta $<$ dadi
variabel relasi rong-papan~$L$, banjur kabeh dikuwantifikasi
universal. Ukara iki valid ing logika orde kapindho murni
yen lan mung yen ukara asal valid, yen lan mung yen
$\Th{PA^2} \Entails !A$, yen lan mung yen
$\Sat{N}{!A}$. Mula yen ana sistem bukti kang sahih
lan jangkep kanggo logika orde kapindho, kita bisa
nggunakake kanggo netepake enumerasi komputabel
$f\colon \Nat \to \Sent[L_A]$ saka !!{sentence}s kang
bener ing~$\Struct{N}$. Fungsi iki bakal bisa
direpresentasikake ing~$\Th{Q}$ dening sawijining
formula orde kapisan~$!B_f(x, y)$. Banjur
!!{formula}~$\lexists[x][!B_f(x, y)]$ bakal netepake
himpunan !!{sentence}s orde kapisan kang bener
ing~$\Struct{N}$, kosok balen karo Teorema Tarski.
\end{proof}



\end{document}
